% Generated by Listas/make_statements.py from pratica-lecture-04.md, item 9. Do not edit.
\textbf{Source.} Tadelis (2013), \emph{Game Theory: An Introduction}, Exercise 8.15 ``Time Inconsistency'', p.~174 {[}PDF 193{]}. Reproduced verbatim. \emph{Note (not in the book):} the example of Section 8.3.4 is on pp.~166--169; its printed first-period choice \(x_1 = \frac{2}{3}K\) (p.~169) does not follow from the objective printed there (the maximum is at \(\frac{K}{2}\)); see \texttt{Notas/L03/L03-08-time-inconsistency.md} §2.2. Answer in the source: none in the book (the preface, p.~xiii, says about half are solved on Princeton UP's website; not checked).



\begin{quote}

\textbf{Time Inconsistency:} Consider the three-period example of a player with hyperbolic discounting described in Section 8.3.4 with utility \(\ln(x)\) in each of the three periods, with initial budget \(K\) and with discount factors \(\delta = 1\) and \(\beta = \frac{1}{2}\).



\begin{enumerate}

\def\labelenumi{\alph{enumi}.}

\item

  Solve the optimal plan of action of a naive player 1 who does not take into account how his future self, player 2, will alter the plan. What is player 1's optimal plan \(x_1^*\), \(x_2^*\), and \(x_3^*\) as a function of \(K\)?

\item

  Let \(K_2\) be the amount left from the solution to (a) after player 1 consumes his planned choice of \(x_1^*\). Given \(K_2\), what is the optimal plan of player 2? In what way does it differ from the optimal plan set out by player 1?

\end{enumerate}

\end{quote}



\textbf{Referred to in the statement.} The three-period example of Section 8.3.4 (Tadelis pp.~166 and 168 {[}PDF p.~185 and 187{]}) (reproduced from the book):



\begin{quote}

Consider a player with \(u(x) = \ln(x)\) who needs to allocate a fixed budget \(K\) across three periods. Since he will not waste any of his budget, it follows that \(x_1 = K - x_2 - x_3\) {[}\ldots{]}



To model this idea consider a different discounting rule known as \emph{hyperbolic discounting}. The player uses the discount rate \(\delta\) for the future periods as in exponential discounting, but he uses an \emph{additional} discount factor \(\beta \in (0, 1)\) to discount all of the future compared to present consumption. Using our three-period example, the player's modified discounted present-value problem is given by



\[\max_{x_2, x_3} v(K - x_2 - x_3, x_2, x_3) = \ln(K - x_2 - x_3) + \beta\delta \ln(x_2) + \beta\delta^2 \ln(x_3). \tag{8.9}\]



Simply put, when the player is looking toward the future, the discount factor he uses between periods \(t = 1\) and \(t = 2\) is stronger than the one he uses between periods \(t = 2\) and \(t = 3\). In particular, at any period and looking forward, the player uses \(\beta\delta < \delta\) to discount between the current period and the next, while he uses only \(\delta\) to discount between any two consecutive future periods.\(^4\) As we will now see, hyperbolic discounting will cause problems of self-control, in that a player will plan to do one thing but later choose to revise his plan.



Imagine that the player is very well aware of his self-control problem---at period \(t = 1\) he discounts between periods \(t = 2\) and \(t = 3\) at the rate of \(\delta\), but he knows that when tomorrow (\(t = 2\)) comes around, he will discount \(t = 3\) at a rate of \(\beta\delta\). In other words, the player at \(t = 1\) (call him player 1) is aware that he is playing a game with his ``future self'' at \(t = 2\) (call him player 2), who will exhibit different preferences over time. In this case a rational, forward-looking player 1 will solve his problem using backward induction; in other words he will solve for the subgame-perfect equilibrium of the implied dynamic game.



\begin{enumerate}

\def\labelenumi{\arabic{enumi}.}

\setcounter{enumi}{3}

\tightlist

\item

  More generally the discounted present value of a player who uses hyperbolic discounting is given by \(u(x_1) + \beta \sum_{t=2}^{T} \delta^{t-1} u(x_t)\).

\end{enumerate}

\end{quote}



\begin{center}\rule{0.5\linewidth}{0.5pt}\end{center}

